% !TEX root = ../main.tex
\begin{ouverture}[rc]{26}{Circuit RC et capteurs capacitifs}{Partie 4 · Ondes et signaux}
{Débit de charges\enspace·\enspace Condensateur\enspace·\enspace Charge et décharge\enspace·\enspace Temps caractéristique\enspace·\enspace Capteurs capacitifs}
\objectif{Relier l'intensité d'un courant électrique au débit de charges.}
\objectif{Identifier des situations où des charges de signes opposés s'accumulent sur des surfaces en regard ; citer des
ordres de grandeur de capacités.}
\objectif{Établir et résoudre l'équation différentielle vérifiée par la tension aux bornes d'un condensateur lors de sa
charge et de sa décharge.}
\objectif{Déterminer le temps caractéristique d'un dipôle RC et expliquer le principe de capteurs capacitifs.}
\objectif{Modéliser la charge d'un condensateur par la méthode d'Euler en Python.}
\prerequis{Intensité, tension, loi d'Ohm, loi des mailles (classes de seconde et de première).}
\prerequis{Lien entre intensité et débit de charges en régime continu (classe de première).}
\prerequis{Équation différentielle $y'=a\,y+b$ et exponentielle (mathématiques, chapitre~20).}
\end{ouverture}
\essentiel
\essSection{1}{Intensité et débit de charges}
\begin{definition}[Intensité d'un courant variable]
L'intensité d'un courant est le débit de charges électriques à travers une section du conducteur. En régime variable,
si la charge $\mathrm{d}q$ traverse la section pendant la durée $\mathrm{d}t$ :
\[i(t)=\frac{\mathrm{d}q}{\mathrm{d}t},\]
avec $i$ en ampères et $q$ en coulombs. En régime continu, on retrouve $I=\dfrac{Q}{\Delta t}$.
\end{definition}
\essSection{2}{Le condensateur}
\begin{definition}[Comportement capacitif et condensateur]
Un \textbf{condensateur} est formé de deux surfaces conductrices en regard, les \textbf{armatures}, séparées par un isolant.
Branché sur un générateur, il accumule des charges opposées, $+q$ et $-q$, sur ses armatures : c'est un comportement
\textbf{capacitif}. La charge $q$ de l'armature positive est proportionnelle à la tension $u_C$ entre les armatures :
\[q=C\,u_C,\]
où $C$ est la \textbf{capacité} du condensateur, en farads (F). On rencontre des capacités de quelques picofarads (capteurs,
circuits rapides) à quelques millifarads, et jusqu'à plusieurs milliers de farads pour les supercondensateurs.
\end{definition}
\schema{26-condensateur}
\begin{propriete}[Influence de la géométrie]
La capacité augmente avec la surface $S$ des armatures en regard, diminue quand leur distance $e$ augmente et dépend de
l'isolant. Pour un condensateur plan, $C=\dfrac{\varepsilon\,S}{e}$ (expression fournie), où $\varepsilon$ caractérise l'isolant :
$\varepsilon_0=\qty{8,85e-12}{F.m^{-1}}$ pour le vide ou l'air. Deux plaques de \qty{1}{m^2} distantes de \qty{1}{mm} dans l'air
n'ont que \qty{8,85}{nF} de capacité.
\end{propriete}
\begin{propriete}[Relation entre intensité et tension]
Le courant qui arrive sur l'armature positive augmente sa charge : en convention récepteur,
\[i=\frac{\mathrm{d}q}{\mathrm{d}t}=C\,\frac{\mathrm{d}u_C}{\mathrm{d}t}.\]
La tension aux bornes d'un condensateur ne peut pas varier brutalement : elle est continue.
\end{propriete}
\essSection{3}{Charge et décharge dans un circuit RC}
\schema{26-circuit}
\begin{savoirfaire}[Établir et résoudre l'équation de la charge]
Le condensateur, initialement déchargé, est relié à l'instant $t=0$ à un générateur idéal de tension $E$ par une
résistance $R$ (position 1). La loi des mailles donne $E=u_R+u_C=R\,i+u_C$, avec $i=C\,\dfrac{\mathrm{d}u_C}{\mathrm{d}t}$ :
\[RC\,\frac{\mathrm{d}u_C}{\mathrm{d}t}+u_C=E\qquad\text{soit}\qquad\frac{\mathrm{d}u_C}{\mathrm{d}t}=-\frac{1}{\tau}\,(u_C-E),\quad\tau=RC.\]
Les solutions s'écrivent $u_C=E+K\,\mathrm{e}^{-t/\tau}$ ; la condition initiale $u_C(0)=0$ impose $K=-E$ :
\[u_C(t)=E\left(1-\mathrm{e}^{-t/\tau}\right).\]
\end{savoirfaire}
\begin{propriete}[Décharge]
En position 2, le condensateur chargé sous la tension $E$ se décharge dans la résistance : $RC\,\dfrac{\mathrm{d}u_C}{\mathrm{d}t}+u_C=0$,
d'où $u_C(t)=E\,\mathrm{e}^{-t/\tau}$. L'intensité, $i=C\,\dfrac{\mathrm{d}u_C}{\mathrm{d}t}=-\dfrac{E}{R}\,\mathrm{e}^{-t/\tau}$, est négative : le courant
circule en sens inverse de celui de la charge.
\end{propriete}
\begin{propriete}[Temps caractéristique]
Le \textbf{temps caractéristique} $\tau=RC$ (en secondes si $R$ est en ohms et $C$ en farads) fixe la durée du régime
transitoire. À $t=\tau$, la charge a atteint 63~\% de sa valeur finale ; au bout de $5\tau$, plus de 99~\% : le régime
permanent est pratiquement atteint. Graphiquement, la tangente à l'origine coupe l'asymptote à $t=\tau$.
\end{propriete}
\schema{26-charge}
\begin{exemple}
Avec $R=\qty{10}{k\ohm}$ et $C=\qty{100}{\micro F}$, $\tau=\num{10e3}\times\num{100e-6}=\qty{1,0}{s}$ : la charge est pratiquement terminée
au bout de \qty{5}{s}.
\end{exemple}
\essSection{4}{Les capteurs capacitifs}
\begin{propriete}[Principe d'un capteur capacitif]
Dans un capteur capacitif, la grandeur à mesurer modifie la capacité d'un condensateur : la distance entre armatures
(capteurs de pression, microphones, accéléromètres), la surface en regard (capteurs de position) ou l'isolant (capteurs
d'humidité, de niveau d'une cuve). Un circuit électronique mesure la capacité, par exemple à travers la durée de charge
d'un circuit RC. Dans un écran tactile, le doigt, conducteur, forme avec une électrode un condensateur supplémentaire ;
dans un smartphone, l'accéléromètre est un micro-condensateur dont une armature mobile se déplace sous l'effet de
l'accélération.
\end{propriete}
\begin{codecote}[0.64]
\begin{pythoncol}
from math import log
E, R, C = 5.0, 10e3, 100e-6      # V, ohm, F
tau, dt = R * C, 0.001           # s
u, t = 0.0, 0.0
while u < 0.95 * E:              # 95 % de E
    u = u + (E - u) / tau * dt
    t = t + dt
print(round(t, 3), round(-tau * log(0.05), 3))
\end{pythoncol}
\begin{consolecol}
2.995 2.996
\end{consolecol}
\cote
\fontsize{9.5}{12.5}\selectfont
\capacite{Capacité numérique}\par\smallskip
Le programme applique la méthode d'Euler à l'équation de la charge, $\dfrac{\mathrm{d}u_C}{\mathrm{d}t}=\dfrac{E-u_C}{\tau}$, jusqu'à ce que
la tension atteigne 95~\% de $E$. Il affiche la durée obtenue et la valeur exacte $-\tau\ln0{,}05\approx3\tau$ : les deux
concordent à la milliseconde près.
\end{codecote}
\begin{attention}
Dans $\tau=RC$, les unités doivent être cohérentes : un kiloohm vaut $10^3$~\unit{\ohm}, un microfarad $10^{-6}$~F. La tension aux
bornes du condensateur est continue, mais l'intensité peut varier brutalement, par exemple à la fermeture de
l'interrupteur, où elle passe de 0 à $E/R$.
\end{attention}
\begin{experience}[Déterminer le temps caractéristique d'un dipôle RC]
Un microcontrôleur commande la charge d'un condensateur à travers une résistance et mesure la tension à ses bornes à
intervalles réguliers ; une carte d'acquisition ou un oscilloscope alimenté par un signal créneau convient aussi. On
détermine $\tau$ par la tangente à l'origine ou par la date à laquelle $u_C=0{,}63\,E$, et on le compare à $RC$. Deux
feuilles d'aluminium séparées par du papier forment un condensateur dont on observe l'augmentation de capacité quand
on presse les feuilles l'une contre l'autre.
\end{experience}
\begin{rediger}
Orienter le circuit, flécher les tensions en convention récepteur, puis écrire la loi des mailles. Remplacer $u_R$ par $R\,i$
et $i$ par $C\,\dfrac{\mathrm{d}u_C}{\mathrm{d}t}$ ; mettre l'équation sous la forme $\dfrac{\mathrm{d}u_C}{\mathrm{d}t}+\dfrac{u_C}{\tau}=\dfrac{E}{\tau}$ avant de
la résoudre, et utiliser la condition initiale.
\end{rediger}
\exercices

\rubrique{Découverte}
\begin{blocLivre}{ex-26-1}
\begin{exercice}[1]{Un débit de charges}
\begin{donnees}
Charge élémentaire : $e=\qty{1,60e-19}{C}$.
\end{donnees}
\begin{questions}
\item Un courant d'intensité constante \qty{2,0}{A} circule pendant \qty{30}{s}. Quelle charge traverse une section du fil ?
\item Combien d'électrons cela représente-t-il ?
\item La charge d'une armature varie selon $q(t)=0{,}5\,t^2$, avec $q$ en millicoulombs et $t$ en secondes. Exprimer
l'intensité du courant en fonction du temps et la calculer à $t=\qty{4}{s}$.
\item Pourquoi la définition $I=Q/\Delta t$ ne convient-elle pas dans ce dernier cas ?
\end{questions}
\end{exercice}
\end{blocLivre}
\begin{blocLivre}{ex-26-2}
\begin{exercice}[1]{Des comportements capacitifs}
\begin{questions}
\item Qu'appelle-t-on comportement capacitif ?
\item Avant un orage, la base des nuages se charge négativement et le sol positivement. Pourquoi peut-on parler de
condensateur ? Quel est l'isolant ?
\item La membrane d'un neurone sépare des ions positifs à l'extérieur et négatifs à l'intérieur. Identifier les
« armatures » et l'isolant.
\item Le doigt posé sur un écran tactile forme un condensateur avec une électrode de l'écran. Quel est l'isolant ?
\end{questions}
\end{exercice}
\end{blocLivre}
\begin{blocLivre}{ex-26-3}
\begin{exercice}[1]{Charge et tension}
Un condensateur de \qty{470}{\micro F} est chargé sous \qty{12}{V}.
\begin{donnees}
$e=\qty{1,60e-19}{C}$.
\end{donnees}
\begin{questions}
\item Calculer la charge de son armature positive.
\item Combien d'électrons ont quitté cette armature ?
\item Même question pour un condensateur de \qty{4,7}{nF} chargé sous la même tension. Comparer.
\item Que devient la charge si la tension double ?
\end{questions}
\end{exercice}
\end{blocLivre}
\begin{blocLivre}{ex-26-4}
\begin{exercice}[1]{Lire une courbe de charge}
On ferme à l'instant $t=0$ l'interrupteur d'un circuit RC série alimenté par un générateur de tension $E$ ; la résistance
vaut $R=\qty{1,0}{k\ohm}$.
\schema{26-ex04}
\begin{questions}
\item Lire la valeur de $E$.
\item Déterminer le temps caractéristique par deux méthodes.
\item En déduire la capacité du condensateur.
\item Au bout de combien de temps peut-on considérer le condensateur comme chargé ?
\end{questions}
\end{exercice}
\end{blocLivre}
\begin{blocLivre}{ex-26-5}
\begin{exercice}[1]{Le flash d'un appareil photo}
Le condensateur d'un flash, de capacité \qty{150}{\micro F}, est chargé sous \qty{300}{V}, puis se décharge en \qty{1}{ms} environ dans la
lampe.
\begin{donnees}
Énergie stockée par un condensateur : $\frac12\,C\,u_C^2$.
\end{donnees}
\begin{questions}
\item Calculer la charge stockée.
\item Calculer l'intensité moyenne du courant pendant la décharge.
\item Calculer l'énergie libérée dans la lampe.
\item Pourquoi utilise-t-on un condensateur plutôt que d'alimenter directement la lampe par la pile ?
\end{questions}
\end{exercice}
\end{blocLivre}
\rubrique{Entraînement}
\begin{blocLivre}{ex-26-6}
\begin{exercice}[2]{Établir l'équation de la décharge}
Un condensateur de capacité $C$, chargé sous la tension $E$, est relié à $t=0$ à une résistance $R$.
\begin{questions}
\item Faire un schéma, orienter le circuit et écrire la loi des mailles.
\item En déduire l'équation différentielle vérifiée par $u_C$.
\item Vérifier que $u_C(t)=E\,\mathrm{e}^{-t/\tau}$ en est solution et satisfait la condition initiale.
\item Quelle fraction de $E$ reste-t-il à $t=2\tau$ ? À quelle date la tension vaut-elle $E/2$ ? Application : $E=\qty{9,0}{V}$.
\end{questions}
\end{exercice}
\end{blocLivre}
\begin{blocLivre}{ex-26-7}
\begin{exercice}[2]{La minuterie d'un escalier}
Une minuterie maintient la lumière allumée tant que la tension aux bornes d'un condensateur de \qty{1000}{\micro F}, chargé sous
\qty{12}{V} à l'appui sur le bouton, reste supérieure à \qty{2,0}{V}. Le condensateur se décharge dans une résistance $R$.
\begin{questions}
\item Exprimer la date à laquelle la lumière s'éteint en fonction de $\tau$.
\item Quelle valeur de $\tau$ donne une durée d'éclairage de \qty{60}{s} ?
\item En déduire la valeur de $R$.
\item Comment allonger la durée d'éclairage sans changer de condensateur ?
\end{questions}
\end{exercice}
\end{blocLivre}
\begin{blocLivre}{ex-26-8}
\begin{exercice}[2]{Modéliser la charge en Python}
On reprend le programme de l'essentiel.
\begin{questions}
\item \capacite{Capacité numérique}Justifier la ligne \code{u = u + (E - u) / tau * dt}.
\item Que se passe-t-il si l'on choisit \code{dt = 0.5} ? Pourquoi faut-il un pas petit devant $\tau$ ?
\item Modifier le programme pour qu'il calcule la durée nécessaire pour atteindre 99~\% de $E$. Comparer à la valeur
exacte.
\item Écrire les lignes qui enregistrent et tracent $u_C$ en fonction du temps.
\end{questions}
\end{exercice}
\end{blocLivre}
\begin{blocLivre}{ex-26-9}
\begin{exercice}[2]{La géométrie d'un condensateur}
Deux plaques métalliques carrées de \qty{20}{cm} de côté sont placées face à face, à \qty{2,0}{mm} l'une de l'autre, dans l'air.
\begin{donnees}
$C=\dfrac{\varepsilon\,S}{e}$ ; dans l'air, $\varepsilon=\varepsilon_0=\qty{8,85e-12}{F.m^{-1}}$ ; avec un plastique, $\varepsilon=3\,\varepsilon_0$.
\end{donnees}
\begin{questions}
\item Calculer la capacité de ce condensateur.
\item On glisse entre les plaques une feuille de plastique de même épaisseur. Calculer la nouvelle capacité.
\item Que devient la capacité dans l'air si l'on rapproche les plaques à \qty{1,0}{mm} ?
\item Quelle surface faudrait-il pour obtenir \qty{1}{\micro F} dans l'air avec un écart de \qty{2,0}{mm} ? Comment fabrique-t-on
alors des condensateurs de forte capacité ?
\end{questions}
\end{exercice}
\end{blocLivre}
\begin{blocLivre}{ex-26-10}
\begin{exercice}[2]{Le défibrillateur}
Un défibrillateur charge un condensateur de \qty{150}{\micro F} sous \qty{2000}{V}, puis le décharge à travers le thorax du patient,
assimilé à une résistance de \qty{50}{\ohm}.
\begin{donnees}
Énergie stockée par un condensateur : $\frac12\,C\,u_C^2$.
\end{donnees}
\begin{questions}
\item Calculer la charge et l'énergie stockées.
\item Calculer le temps caractéristique de la décharge et l'intensité du courant au début de la décharge.
\item Calculer la tension aux bornes du condensateur au bout de \qty{10}{ms}.
\item Pourquoi la décharge est-elle réalisée en quelques millisecondes seulement ?
\end{questions}
\end{exercice}
\end{blocLivre}
\begin{blocLivre}{ex-26-11}
\begin{exercice}[2]{Le signal créneau}
Un générateur délivre une tension créneau qui passe périodiquement de 0 à \qty{5}{V}, puis de \qty{5}{V} à 0 ; il alimente un
circuit RC avec $R=\qty{1,0}{k\ohm}$ et $C=\qty{100}{nF}$. On observe $u_C$ à l'oscilloscope.
\begin{questions}
\item Calculer le temps caractéristique du circuit.
\item Pendant chaque demi-période, le condensateur se charge ou se décharge. Quelle durée minimale faut-il pour que la
charge soit pratiquement complète ?
\item En déduire la fréquence maximale du créneau qui permet d'observer des charges complètes.
\item Quelle tension le condensateur atteint-il, en partant de 0, à la fin de la première demi-période si la fréquence vaut
\qty{5}{kHz} ?
\end{questions}
\end{exercice}
\end{blocLivre}
\rubrique{Objectif bac}
\begin{blocLivre}{ex-26-12}
\begin{exercice}[3]{L'écran tactile}
\begin{docu}{1}{Détecter un doigt}
Sous la vitre d'un écran tactile, chaque électrode forme avec son environnement un condensateur d'environ \qty{10}{pF}.
Quand un doigt touche la vitre au-dessus d'elle, sa capacité augmente d'environ \qty{1}{pF}.
\end{docu}
\begin{docu}{2}{La mesure}
Le circuit charge périodiquement l'électrode à travers une résistance de \qty{1,0}{M\ohm} et mesure la durée nécessaire pour
que sa tension atteigne la moitié de la tension d'alimentation.
\end{docu}
\begin{questions}
\item Calculer le temps caractéristique de charge sans doigt, puis avec le doigt.
\item Montrer que la durée nécessaire pour atteindre la moitié de la tension finale vaut $\tau\ln2$.
\item \initiative Calculer la différence de durée que le circuit doit détecter. Commenter.
\item Pourquoi un écran capacitif ne réagit-il pas au contact d'un stylo en plastique ou d'un doigt ganté ?
\end{questions}
\end{exercice}
\end{blocLivre}
\begin{blocLivre}{ex-26-13}
\begin{exercice}[3]{L'accéléromètre du smartphone}
\begin{docu}{1}{Un micro-condensateur}
Dans l'accéléromètre d'un smartphone, une masse mobile de \qty{2,0e-9}{kg}, suspendue à un ressort de raideur
$k=\qty{1,0}{N.m^{-1}}$, porte une armature d'un condensateur plan. L'autre armature est fixe, à $e=\qty{2,0}{\micro m}$ ; la
surface en regard vaut \qty{1,0e-7}{m^2}.
\end{docu}
\begin{docu}{2}{Le principe}
Sous l'effet d'une accélération $a$, la masse se déplace de $x=\dfrac{m\,a}{k}$, ce qui modifie l'écart entre les armatures. Pour
$x\ll e$, la variation relative de capacité vaut $\dfrac{\Delta C}{C}\approx\dfrac{x}{e}$.
\end{docu}
\begin{donnees}
$C=\dfrac{\varepsilon_0\,S}{e}$ avec $\varepsilon_0=\qty{8,85e-12}{F.m^{-1}}$ ; $g=\qty{9,81}{m.s^{-2}}$.
\end{donnees}
\begin{questions}
\item Calculer la capacité au repos.
\item Calculer le déplacement de la masse pour une accélération égale à $g$.
\item \initiative Calculer la variation de capacité correspondante. Commenter son ordre de grandeur.
\item Comment le smartphone sait-il s'il est tenu verticalement ou horizontalement ?
\end{questions}
\end{exercice}
\end{blocLivre}
\begin{blocLivre}{ex-26-14}
\begin{exercice}[3]{Le capteur de niveau d'une cuve}
\begin{docu}{1}{Une sonde capacitive}
Une sonde verticale, plongée dans une cuve de \qty{2,5}{m} de hauteur, forme un condensateur dont l'isolant est en partie de
l'air et en partie le liquide. Sa capacité varie linéairement avec la hauteur $h$ de liquide.
\end{docu}
\begin{docu}{2}{L'étalonnage}
La sonde est chargée à travers une résistance fixe. Le temps caractéristique vaut \qty{50}{\micro s} quand la cuve est vide et
\qty{120}{\micro s} quand elle est pleine. Une mesure donne \qty{78}{\micro s}.
\end{docu}
\begin{questions}
\item Pourquoi le temps caractéristique varie-t-il comme la capacité ?
\item Justifier qu'il varie linéairement avec la hauteur de liquide.
\item \initiative Déterminer la hauteur de liquide dans la cuve.
\item Quel temps caractéristique mesurera-t-on lorsque la cuve sera remplie aux trois quarts ?
\end{questions}
\end{exercice}
\end{blocLivre}
\begin{blocLivre}{ex-26-15}
\begin{exercice}[3]{Le stimulateur cardiaque}
\begin{docu}{1}{Un rythme imposé}
Dans un modèle simple de stimulateur cardiaque, un condensateur de \qty{1,0}{\micro F}, chargé sous \qty{2,8}{V}, se décharge à travers
une résistance $R$. Lorsque sa tension descend à $E/\mathrm{e}$, soit environ \qty{1,03}{V}, le circuit envoie une impulsion au cœur et
recharge aussitôt le condensateur.
\end{docu}
\begin{docu}{2}{Le réglage}
Le médecin règle le stimulateur sur 70 battements par minute.
\end{docu}
\begin{questions}
\item Montrer que la durée entre deux impulsions est égale au temps caractéristique $\tau$.
\item Calculer la durée entre deux battements.
\item \initiative En déduire la valeur de la résistance.
\item Quelle valeur de $R$ faut-il pour un rythme de 90 battements par minute pendant un effort ?
\end{questions}
\end{exercice}
\end{blocLivre}
